% Part: lambda-calculus
% Chapter: syntax
% Section: abbreviated-syntax

\documentclass[../../../include/open-logic-section]{subfiles}

\begin{document}

\olfileid{lam}{syn}{abb}
\olsection{સંક્ષિપ્ત વાક્યરચના}

\olref[trm]{defn:term} મુજબ વ્યાખ્યાયિત પદો લખવામાં ક્યારેક
ભારે પડે છે; તેથી વધુ સંક્ષિપ્ત વાક્યરચના ઉપયોગી છે. આ
સંક્ષિપ્ત સંકેતનમાં પણ પદોને અનન્ય રીતે વાંચી શકાય તેની
કાળજી લેવી જોઈએ. વ્યાપકપણે વપરાતું એક સ્વરૂપ
\emph{સંક્ષિપ્ત પદો} કહેવાય છે અને તે આ પ્રમાણે છે.

\begin{enumerate}
\item કૌંસ છોડવામાં આવે ત્યારે અનુપ્રયોગ ડાબેથી જમણે થાય છે.
  દાખલા તરીકે, $M$, $N$, $P$ અને~$Q$ પદો હોય તો
  $MNPQ$ એ~$(((MN)P)Q)$નું સંક્ષિપ્ત સ્વરૂપ છે.
\item ફરી, કૌંસ છોડવામાં આવે ત્યારે લૅમ્બડા અમૂર્તીકરણનો
  વ્યાપ શક્ય તેટલો મોટો લેવાય છે. દાખલા તરીકે,
  $\lambd[x][MNP]$ને~$(\lambd[x][MNP])$ એમ વાંચવું.
\item એક લૅમ્બડાથી અનેક ચલોનું અમૂર્તીકરણ કરી શકાય છે.
  દાખલા તરીકે, $\lambd[xyz][M]$ એ
  $\lambd[x][\lambd[y][\lambd[z][M]]]$નું સંક્ષિપ્ત સ્વરૂપ છે.
\end{enumerate}

દાખલા તરીકે,
\[
\lambd[xy][xxyx \lambd[z][xz]]
\]
એ નીચેનું સંક્ષિપ્ત સ્વરૂપ છે:
\[
(\lambd[x][(\lambd[y][((((xx)y)x)(\lambd[z][(xz)]))])]).
\]

\begin{prob}
સંક્ષિપ્ત પદ~$\lambd[g][(\lambd[x][g (x x)])
  \lambd[x][g (x x)]]$ને પૂર્ણ સ્વરૂપે લખો.
\end{prob}

\end{document}
